\documentclass[10pt,a4paper]{article}

\usepackage[T1]{fontenc}
\usepackage[utf8]{inputenc}
\usepackage{amsmath,amssymb,amsthm}
\usepackage{mathtools}
\usepackage{geometry}
\usepackage{graphicx}
\usepackage{hyperref}
\usepackage{booktabs}
\usepackage{algorithm}
\usepackage{algpseudocode}
\usepackage{xcolor}
\usepackage{cite}

\geometry{margin=2cm}

\newtheorem{theorem}{Théorème}
\newtheorem{lemma}[theorem]{Lemme}
\newtheorem{proposition}[theorem]{Proposition}
\newtheorem{corollary}[theorem]{Corollaire}
\newtheorem{definition}{Définition}
\newtheorem{remark}{Remarque}

\title{\textbf{B-MoE : Mixture of Experts à Routage par Blocs\\
et Spécialisation Guidée des Experts}\\[0.5em]
\large Technical Report v3 — Proposition, validation empirique et raffinement architectural}

\author{[Sébastien Viou]}

\date{Juin 2026}

\begin{document}

\maketitle

\begin{abstract}
Nous proposons \textbf{B-MoE} (\textit{Bloc-routed Mixture of Experts}), une architecture étendant les Mixture of Experts (MoE) selon trois axes complémentaires : (1) un \textbf{routage par blocs} opérant tous les $Z$ couches, réduisant le nombre de décisions discrètes et stabilisant le gradient ; (2) un \textbf{mécanisme de scoring par attention inter-experts} à chaque point de routage, permettant une décision informée par la sémantique collective des experts actifs ; (3) une \textbf{initialisation guidée par domaine} (\textit{pre-specialization}), inspirée de la modularité fonctionnelle corticale, accélérant la convergence et garantissant la diversité des représentations. Par rapport au travail le plus proche, RMoE \cite{qiu2025rmoe}, qui introduit un GRU inter-couches pour propager les décisions de routage, B-MoE se distingue sur trois points : le routage opère sur des blocs de $Z$ couches et non couche-par-couche, le mécanisme de scoring repose sur une attention inter-experts sur les sorties agrégées du bloc plutôt que sur un état récurrent scalaire, et la spécialisation est initialisée \textit{a priori} avant l'entraînement joint plutôt qu'émergente. Nous établissons les propriétés théoriques principales de ces trois mécanismes et discutons leur intégration dans un pipeline d'entraînement unifié. Au-delà de la proposition initiale, cette version~3 ajoute une \textbf{validation empirique} sur un banc d'essai synthétique contrôlé : la commutation d'experts entre blocs est \emph{fonctionnellement nécessaire} à la composition, mais des experts liés à une profondeur fixe \emph{ne généralisent pas} aux compositions inédites (zéro-shot). Nous proposons alors un \textbf{raffinement architectural} qui lève cette limite --- une \emph{boucle d'experts partagés à ré-ancrage par tokens} (récurrence à la Universal Transformer~\cite{dehghani2019universal} dotée du routage par experts) --- ainsi qu'un \textbf{routage autonome} par démonstrations \emph{few-shot} et \emph{critique de suffisance} (\textit{propose-and-verify}), qui résout la généralisation compositionnelle zéro-shot (précision de $1.00$ sur des compositions jamais vues, paires et triples). Nous précisons enfin la portée (banc synthétique) et la voie de passage à l'échelle.
\end{abstract}

\section{Introduction}

Les architectures Mixture of Experts \cite{shazeer2017outrageously} permettent d'augmenter massivement la capacité paramétrique d'un modèle sans augmentation proportionnelle du coût computationnel. Dans les formulations actuelles (Switch Transformer \cite{fedus2022switch}, Mixtral \cite{jiang2024mixtral}), le routage est effectué indépendamment à chaque couche, sans mémoire des décisions antérieures.

Cette indépendance couche-par-couche présente trois limitations fondamentales :
\begin{enumerate}
  \item Le router à la couche $\ell$ ignore quels experts ont été activés aux couches $\ell-1, \ldots, \ell-k$, perdant une information de contexte précieuse sur la trajectoire d'activation.
  \item La multiplicité des décisions de routage introduit du bruit dans le gradient sur les réseaux profonds, fragilisant la convergence des routers.
  \item Sans contrainte de spécialisation, les experts tendent vers la redondance (\textit{expert collapse}), réduisant l'efficacité paramétrique effective du modèle.
\end{enumerate}

RMoE \cite{qiu2025rmoe} a récemment adressé la première limitation en introduisant un GRU pour propager l'état de routage couche-par-couche. B-MoE adopte une stratégie différente et complémentaire : au lieu de propager un état récurrent à chaque couche, nous groupons les couches en \textit{blocs} de taille $Z$, effectuons le routage à la granularité du bloc via une attention inter-experts, et initialisons les experts par domaine avant l'entraînement joint.

\section{Travaux connexes}

\paragraph{Routage inter-couches.}
Le routage standard dans les MoE \cite{shazeer2017outrageously,fedus2022switch} est indépendant à chaque couche. RMoE \cite{qiu2025rmoe} est le travail le plus directement connexe : il introduit un GRU léger entre les routers de couches consécutives, propageant l'information de routage couche par couche. Notre approche diffère en deux points : (i) le routage opère à la granularité du \textit{bloc} ($Z$ couches), réduisant le nombre de décisions de $L$ à $B = L/Z$ ; (ii) le mécanisme de scoring repose sur une attention sur les sorties agrégées du bloc, pas sur un état GRU scalaire. Les deux approches sont orthogonales et potentiellement composables.

\paragraph{Spécialisation des experts.}
Branch-Train-Merge \cite{li2022btm} entraîne des experts séparément sur des domaines distincts puis les fusionne — spécialisation explicite mais post-hoc, sans entraînement joint. UpIT \cite{upit2025} assigne des données spécialisées à chaque expert lors d'un upcycling depuis un dense, avec une loss de classification auxiliaire. Expert Divergence Learning \cite{edl2026} guide la spécialisation en décomposant le corpus d'entraînement par domaine. B-MoE se distingue en initialisant les experts \textit{avant} l'entraînement joint (pré-spécialisation \textit{a priori}), ce qui oriente le paysage de loss dès le début de l'optimisation plutôt qu'en contraignant le routage en cours d'entraînement.

\paragraph{Routage par profondeur et blocs.}
Mixture-of-Depths \cite{raposo2024mod} alloue dynamiquement le compute par token selon la profondeur, mais ne fixe pas d'expert par bloc. Les Universal Transformers \cite{dehghani2019universal} réutilisent les mêmes poids sur plusieurs passes, ce qui est adjacent à notre notion de bloc mais sans routing différencié. À notre connaissance, aucun travail ne combine routage par blocs, scoring par attention inter-experts et pré-spécialisation a priori dans un cadre unifié.

\section{Formalisation du cadre}

\subsection{Notations}

Soit un transformer de profondeur $L$ composé de $N$ experts par position de routage. Notons :
\begin{itemize}
  \item $\mathbf{h}_\ell^{(t)} \in \mathbb{R}^d$ : le vecteur du \textit{residual stream} au token $t$ après la couche $\ell$
  \item $E_i^{(\ell)} : \mathbb{R}^d \to \mathbb{R}^d$ : la fonction de l'expert $i$ à la couche $\ell$
  \item $Z \in \mathbb{N}^*$ : la taille d'un bloc (nombre de couches par décision de routage)
  \item $B = \lfloor L/Z \rfloor$ : le nombre total de blocs
  \item $\sigma_b \in \{1, \ldots, N\}$ : l'indice de l'expert actif dans le bloc $b$
\end{itemize}

\subsection{Structure en blocs}

\begin{definition}[Bloc B-MoE]
Le bloc $b \in \{1, \ldots, B\}$ correspond aux couches $\ell \in \{(b-1)Z+1, \ldots, bZ\}$. Au sein d'un bloc, un unique expert $\sigma_b$ est actif pour l'ensemble des $Z$ couches. Le residual stream se propage selon :
\begin{equation}
  \mathbf{h}_\ell^{(t)} = \mathbf{h}_{\ell-1}^{(t)} + E_{\sigma_b}^{(\ell)}\!\left(\mathbf{h}_{\ell-1}^{(t)}\right), \quad \ell \in \text{bloc } b
\end{equation}
La représentation agrégée à la sortie du bloc $b$ est :
\begin{equation}
  \mathbf{o}_b^{(t)} = \mathbf{h}_{bZ}^{(t)}
\end{equation}
\end{definition}

\begin{remark}[Comparaison avec RMoE]
Dans RMoE \cite{qiu2025rmoe}, le GRU propage un état caché $\mathbf{r}_\ell \in \mathbb{R}^{d_r}$ couche par couche : $\mathbf{r}_\ell = \mathrm{GRU}(\mathbf{r}_{\ell-1}, \mathbf{h}_\ell)$. Le router à la couche $\ell$ conditionne sa décision sur $\mathbf{r}_{\ell-1}$. B-MoE supprime la granularité couche-par-couche : la décision de routage est prise une seule fois par bloc, sur la base de la sortie agrégée $\mathbf{o}_{b-1}$, réduisant le nombre de décisions discrètes de $L$ à $B$.
\end{remark}

\section{Routage inter-blocs par attention}

\subsection{Limitation du scoring scalaire}

Dans les MoE standards, le router est un MLP $r : \mathbb{R}^d \to \mathbb{R}^N$ produisant un score scalaire par expert. Dans RMoE, le GRU enrichit ce score d'un état récurrent mais reste un vecteur de dimension $d_r$ sans modéliser explicitement les relations entre experts. B-MoE introduit un scoring par attention sur les représentations de chaque expert.

\subsection{Attention inter-experts}

Soit $\mathbf{O}_b \in \mathbb{R}^{N \times d}$ la matrice des représentations agrégées de chaque expert au bloc $b$ (voir section \ref{sec:estimation} pour les experts non actifs).

\begin{definition}[Router par attention inter-experts]
Le score de routage pour le bloc $b+1$ est calculé par :
\begin{equation}
  \mathbf{s}_{b+1} = \mathrm{softmax}\!\left(\frac{\mathbf{Q}_b \mathbf{K}_b^\top}{\sqrt{d_k}}\right)\mathbf{V}_b \in \mathbb{R}^{N}
\end{equation}
où $\mathbf{Q}_b = \mathbf{O}_b W_Q$, $\mathbf{K}_b = \mathbf{O}_b W_K$, $\mathbf{V}_b = \mathbf{O}_b W_V$, avec $W_Q, W_K \in \mathbb{R}^{d \times d_k}$ et $W_V \in \mathbb{R}^{d \times 1}$. L'expert sélectionné est :
\begin{equation}
  \sigma_{b+1} = \arg\max_{i}\, s_{b+1,i}
\end{equation}
\end{definition}

\subsection{Estimation des experts non actifs}
\label{sec:estimation}

Un expert $i \neq \sigma_b$ n'a pas produit de sortie au bloc $b$. Trois stratégies sont envisageables :
\begin{enumerate}
  \item \textbf{Propagation du bloc précédent} : $\mathbf{o}_{b,i} \approx \mathbf{o}_{b-1,i}$ — approximation simple, coût nul.
  \item \textbf{Forward partiel} : évaluer $E_i$ sur un sous-échantillon de tokens pour estimer $\mathbf{o}_{b,i}$ — plus précis, surcoût linéaire en $N$.
  \item \textbf{Projection apprise} : $\mathbf{o}_{b,i} = f_i(\mathbf{o}_{b-1,i}, \mathbf{h}_{bZ})$ avec $f_i$ un MLP léger — compromis.
\end{enumerate}

La stratégie (1) est suffisante pour un prototype et constitue l'hypothèse par défaut dans ce rapport. La stratégie (3) est recommandée pour un déploiement à grande échelle.

\begin{proposition}[Différentiabilité du routage]
Le score $\mathbf{s}_{b+1}$ est différentiable par rapport aux paramètres $W_Q, W_K, W_V$ et aux représentations $\mathbf{O}_b$. La sélection $\arg\max$ est rendue différentiable via Gumbel-Softmax \cite{jang2017categorical} ou l'estimateur Straight-Through \cite{bengio2013estimating} lors de l'entraînement.
\end{proposition}

\begin{proof}
$\mathbf{s}_{b+1}$ est une composition de fonctions différentiables en $W_Q, W_K, W_V$ et $\mathbf{O}_b$. Seule l'opération $\arg\max$ est discontinue. En substituant par $\mathrm{softmax}(\mathbf{s}_{b+1}/\tau)$ avec $\tau \to 0$ (Gumbel-Softmax) ou via le gradient straight-through sur la sélection binaire, le graphe de calcul reste différentiable de bout en bout. $\square$
\end{proof}

\section{Pré-spécialisation des experts}

\subsection{Motivation}

Dans le cortex cérébral, la modularité fonctionnelle est partiellement innée : des zones dédiées au langage, à la perception visuelle ou au raisonnement spatial présentent des prédispositions structurelles avant toute expérience \cite{dehaene2005neural}. Cette initialisation guidée accélère la spécialisation post-natale sans l'imposer rigidement.

Contrairement à Branch-Train-Merge \cite{li2022btm} qui spécialise les experts séparément puis les fusionne (sans entraînement joint ultérieur), et à UpIT \cite{upit2025} qui opère en upcycling depuis un dense, B-MoE pré-spécialise les experts \textit{avant} l'entraînement joint, orientant le paysage de loss dès l'initialisation.

\subsection{Procédure}

\begin{definition}[Corpus de pré-spécialisation]
Soit $\mathcal{D} = \{\mathcal{D}_1, \ldots, \mathcal{D}_N\}$ une partition de corpus par domaine. Chaque expert $E_i$ est pré-entraîné sur $\mathcal{D}_i$ pendant $T_{\mathrm{pre}}$ étapes :
\begin{equation}
  \mathcal{L}_{\mathrm{pre}}(E_i) = -\mathbb{E}_{x \sim \mathcal{D}_i}\left[\log p_{E_i}(x)\right]
\end{equation}
\end{definition}

\begin{theorem}[Divergence garantie des experts]
Soient $E_i$ et $E_j$ pré-spécialisés sur des corpus $\mathcal{D}_i$ et $\mathcal{D}_j$ disjoints. Sous des hypothèses de régularité standard (loss localement $\mu$-fortement convexe), les paramètres $\theta_i^*$ et $\theta_j^*$ issus de la minimisation de $\mathcal{L}_{\mathrm{pre}}$ vérifient :
\begin{equation}
  \|\theta_i^* - \theta_j^*\|_2 \geq \delta(\mathcal{D}_i, \mathcal{D}_j) > 0
\end{equation}
où $\delta$ est une fonction croissante de $D_{\mathrm{KL}}(p_{\mathcal{D}_i} \| p_{\mathcal{D}_j})$.
\end{theorem}

\begin{proof}[Esquisse]
Par convergence du SGD, $\theta_i^*$ minimise localement $\mathcal{L}_{\mathrm{pre}}(E_i)$. Si $D_{\mathrm{KL}}(p_{\mathcal{D}_i} \| p_{\mathcal{D}_j}) > 0$, les gradients attendus $\mathbb{E}_{x \sim \mathcal{D}_i}[\nabla \mathcal{L}]$ et $\mathbb{E}_{x \sim \mathcal{D}_j}[\nabla \mathcal{L}]$ diffèrent, induisant des minimiseurs distincts. La borne inférieure découle de la $\mu$-forte convexité locale et de l'écart entre gradients. $\square$
\end{proof}

\subsection{Entraînement joint}

Après pré-spécialisation, l'entraînement joint maintient la diversité via deux termes de régularisation :
\begin{equation}
  \mathcal{L}_{\mathrm{div}} = -\lambda \sum_{i \neq j} D_{\mathrm{KL}}\!\left(p_{\theta_i} \| p_{\theta_j}\right)
\end{equation}
\begin{equation}
  \mathcal{L}_{\mathrm{bal}} = \alpha \sum_{b=1}^{B} \sum_{i=1}^{N} \left(f_{b,i} - \frac{1}{N}\right)^2
\end{equation}
où $f_{b,i}$ est la fraction de tokens routés vers l'expert $i$ au bloc $b$. La loss totale est :
\begin{equation}
  \mathcal{L} = \mathcal{L}_{\mathrm{LM}} + \mathcal{L}_{\mathrm{div}} + \mathcal{L}_{\mathrm{bal}}
\end{equation}

\section{Propriétés théoriques}

\begin{proposition}[Réduction du vanishing gradient sur les routers]
Dans un MoE à routage couche-par-couche, le gradient de la loss par rapport aux paramètres du router à la couche $\ell$ décroît avec la profondeur. Dans B-MoE, le nombre de décisions de routage est réduit à $B = L/Z$, et le gradient intra-bloc se propage sans décision discrète intermédiaire via les connexions résiduelles :
\begin{equation}
  \left\|\frac{\partial \mathcal{L}}{\partial W_{\mathrm{router}}^{(b)}}\right\|_{\mathrm{B\text{-}MoE}} \;\geq\; \left\|\frac{\partial \mathcal{L}}{\partial W_{\mathrm{router}}^{(\ell)}}\right\|_{\mathrm{standard}} \cdot C^{Z-1}
\end{equation}
pour une constante $C > 1$ liée à la norme spectrale des couches du bloc.
\end{proposition}

\begin{proof}
Dans le routage standard, le gradient doit traverser $L$ opérations discrètes. Avec le routage par blocs, il en traverse $B = L/Z$. Le residual stream intra-bloc propage le gradient directement à travers $Z$ couches continues, préservant sa norme via les skip connections. $\square$
\end{proof}

\begin{proposition}[Complexité computationnelle]
Le surcoût du router inter-experts par rapport à un MLP scalaire standard est $\mathcal{O}(N^2 d_k / Z)$ par couche, amorti sur $Z$ couches. Pour $N \leq 16$, $d_k \ll d$ et $Z \geq 4$, ce surcoût est négligeable devant le coût des couches transformer.
\end{proposition}

\section{Validation empirique et raffinement de l'architecture}
\label{sec:empirique}

Nous validons les mécanismes de B-MoE sur un banc d'essai synthétique contrôlé, puis nous identifions une limite de la formulation par blocs liés à la profondeur et proposons un raffinement qui la lève. Toutes les mesures sont sur des données de \emph{test} disjointes de l'entraînement.

\subsection{Protocole}
Nous définissons $N$ \emph{compétences atomiques} comme des applications déterministes sur un vocabulaire partagé (p.~ex. \texttt{math}: $x\mapsto x+5$ ; \texttt{histoire}: $x\mapsto ax+b$ ; \texttt{géographie}: une permutation fixe $\pi$). Un expert est pré-spécialisé par compétence (§5). Une \emph{requête complexe} est une composition, $y = f_k \circ \cdots \circ f_1(x)$ ; la résoudre exige d'enchaîner plusieurs experts.

\subsection{La commutation inter-blocs est fonctionnellement nécessaire}
Sur les compositions \emph{vues} à l'entraînement, nous comparons le routage appris au fait de \emph{forcer un unique expert} sur tous les blocs (ablation).

\begin{center}
\begin{tabular}{lccc}
\toprule
Tâche & Routage appris & Meilleur expert seul & Gain de commutation \\
\midrule
pure (1 compétence) & $0.99$ & $0.98$ & $+0.01$ \\
composée (2--3 compétences) & $0.95$--$0.99$ & $0.02$--$0.26$ & $+0.7$ à $+0.96$ \\
\bottomrule
\end{tabular}
\end{center}

Aucun expert isolé ne résout une requête composée ; seule la commutation d'experts entre blocs y parvient. Ceci valide la thèse centrale (§1).

\subsection{Limite : pas de généralisation compositionnelle zéro-shot}
Lorsque l'on \emph{retient} des compositions entières (entraînement sur certaines paires, test sur une paire et un triple jamais vus), la formulation par blocs \emph{échoue} (précision $\approx 0.01$), et même une recherche \emph{oracle} sur tous les chemins d'experts plafonne à $\approx 0.06$. Le diagnostic est sans appel : il ne s'agit pas d'un défaut de routage mais de \textbf{composabilité}. Parce que la pré-spécialisation entraîne chaque expert sur l'ensemble des blocs, « un bloc de l'expert~$i$ » ne réalise pas « une application réutilisable de la compétence~$i$ » ; de plus, l'expert d'un bloc se \emph{co-adapte} à la distribution produite par son prédécesseur. L'interface inter-bloc n'est donc pas \emph{canonique}, et les compétences ne sont pas réutilisables d'une position à l'autre. Ceci entre en tension avec la nature \emph{dépendante de la profondeur} des experts ($E_i^{(\ell)}$, §3.1) : une compétence liée à une couche n'est pas réutilisable ailleurs.

\subsection{Raffinement : boucle d'experts partagés à ré-ancrage par tokens}
Nous séparons deux rôles aujourd'hui confondus dans la pile de couches. Un \textbf{cœur} (\textit{core}) absorbe \emph{une seule fois} la montée d'abstraction et produit une représentation de travail stationnaire $\mathbf{h}$. Une \textbf{boucle} d'experts \emph{partagés} applique ensuite les compétences à ce niveau stationnaire, en \textbf{ré-ancrant} le flux en espace token entre chaque étape :
\begin{align}
\mathbf{h} &\leftarrow \mathrm{norm}\!\big(\mathbf{h} + E_{\sigma_n}(\mathrm{ctx}(\mathbf{h}))\big), \\
\mathbf{h} &\leftarrow \mathrm{softmax}\!\big(\mathrm{Head}(\mathbf{h})\big)\, \mathbf{W}_{\mathrm{emb}} \qquad \text{(ré-ancrage)} .
\end{align}
Le ré-ancrage garantit que chaque expert reçoit un \emph{token propre}, indépendant de son prédécesseur : chaque expert devient une application réutilisable \texttt{token}$\to$\texttt{token}. C'est une récurrence à la Universal Transformer~\cite{dehghani2019universal} dotée du routage par experts de B-MoE, et cela résout la tension profondeur/abstraction (le cœur gère l'abstraction, la boucle gère l'application des compétences).

\begin{remark}[Résultat empirique]
Avec ré-ancrage, les compositions \emph{jamais vues} --- paire \emph{et} triple --- se composent à une précision de $1.00$ en zéro-shot (contre $\approx 0.07$ sans ré-ancrage). Une recherche sur les chemins retrouve seule la bonne chaîne (rang $1$). Le ré-ancrage \emph{dur} (espace token) est l'ingrédient décisif ; une supervision profonde \emph{molle} ne corrige que les paires.
\end{remark}

\subsection{Routage autonome : démonstrations et critique de suffisance}
Le routeur paramétrique du §4.2 peine à composer de façon autonome : une \emph{démonstration} $(x\!\to\!y)$ ne révèle que l'\emph{application nette}, pas sa \emph{décomposition} en compétences (plusieurs chaînes produisent la même application). Nous communiquons donc la tâche par du \textbf{contenu} --- quelques démonstrations \emph{few-shot} --- sans jamais nommer d'expert. Le modèle \textbf{propose} des chaînes de ses propres primitives, et un \textbf{critique de suffisance} juge « le but est-il atteint ? » en vérifiant chaque chaîne sur les démonstrations ($\text{chaîne}(x_{\text{démo}}) = y_{\text{démo}}$) ; la plus courte chaîne suffisante est appliquée à la requête. Ce schéma \emph{propose-and-verify} lève l'ambiguïté en \emph{composant les primitives et en vérifiant contre le but}.

\begin{remark}[Résultat empirique]
Le modèle \emph{se route lui-même} à partir des seules démonstrations --- sans routeur ni token d'expert --- et atteint une précision de $1.00$ en zéro-shot sur la paire \emph{et} le triple inédits.
\end{remark}

\subsection{Portée et passage à l'échelle}
Ces résultats valident le \emph{mécanisme} sur un banc synthétique : compétences déterministes, critique par correspondance exacte sur les démonstrations, routage par recherche exhaustive ($N^K$). Deux composants doivent être remplacés pour passer à l'échelle, le reste de l'architecture étant transposable : le cœur devient le \emph{backbone} pré-entraîné ; les experts deviennent des adaptateurs (\emph{LoRA}) ou des experts MoE ; le ré-ancrage devient une boucle de \emph{chaîne-de-pensée}~\cite{wei2022cot} (les étapes intermédiaires re-deviennent des tokens) ; le critique exact devient un \textbf{vérificateur appris/sémantique} (« cette réponse est-elle complète ? »)~\cite{cobbe2021verifiers} ; et la recherche exhaustive devient un \textbf{proposeur guidé} (boucle d'agent, \emph{beam} / \emph{best-of-N}) avec arrêt piloté par le critique~\cite{banino2021ponder}. L'\emph{intelligence} du juge de suffisance --- absente d'un petit MLP-jouet --- est attendue comme une propriété d'échelle ; nous distinguons donc le mécanisme (validé en jouet) de cette intelligence (dépendante de l'échelle).

\section{Comparaison avec l'état de l'art}

\begin{center}
\begin{tabular}{lcccc}
\toprule
Architecture & Routage inter-couches & Scoring & Spécialisation & Stabilité gradient \\
\midrule
Switch Transformer \cite{fedus2022switch} & \texttimes & Scalaire & \texttimes & Faible \\
Mixtral \cite{jiang2024mixtral} & \texttimes & Scalaire (top-2) & \texttimes & Faible \\
RMoE \cite{qiu2025rmoe} & GRU couche/couche & Scalaire & \texttimes & Moyenne \\
Branch-Train-Merge \cite{li2022btm} & \texttimes & --- & Post-hoc & N/A \\
UpIT \cite{upit2025} & \texttimes & Scalaire & Upcycling & Moyenne \\
\textbf{B-MoE (ce travail)} & \textbf{Attention/bloc} & \textbf{Inter-experts} & \textbf{A priori} & \textbf{Élevée} \\
\bottomrule
\end{tabular}
\end{center}

\section{Discussion et limitations}

\paragraph{Valeur de $Z$.}
C'est le paramètre le plus sensible. Nous recommandons $Z \in [L/16, L/8]$ comme plage initiale, soit $Z \in [6, 12]$ pour $L = 96$.

\paragraph{Limitation principale : estimation de $\mathbf{O}_b$ pour les experts non actifs.}
La stratégie (1) (propagation du bloc précédent) introduit un biais : si la représentation d'un expert non actif a fortement divergé depuis le bloc précédent, le score de routage peut être mal calibré. Cette limitation est le point théorique le plus vulnérable de B-MoE et constitue une priorité pour la validation empirique. La validation de la section~\ref{sec:empirique} montre toutefois que le verrou \emph{dominant} pour la généralisation compositionnelle n'est pas l'estimation de $\mathbf{O}_b$, mais le caractère \emph{non canonique} de l'interface inter-bloc --- levé par le ré-ancrage par tokens.

\paragraph{Définition des domaines de pré-spécialisation.}
La partition $\{\mathcal{D}_i\}$ requiert une taxonomie préalable qui peut être sous-optimale pour des tâches hybrides. Une extension naturelle est d'apprendre la partition via clustering sur les embeddings du corpus.

\paragraph{Composabilité avec RMoE.}
B-MoE et RMoE sont orthogonaux : le mécanisme GRU de RMoE pourrait être appliqué à l'intérieur des blocs de B-MoE pour propager l'information intra-bloc, tandis que l'attention inter-experts opèrerait au niveau inter-blocs. Cette composition est une direction de recherche future.

\section{Conclusion}

B-MoE propose une architecture MoE combinant trois contributions complémentaires : routage par blocs pour la stabilité du gradient, attention inter-experts pour un scoring sémantiquement informé, et pré-spécialisation guidée pour la diversité et la convergence. La différenciation principale par rapport à RMoE \cite{qiu2025rmoe} réside dans la granularité du routage (bloc vs couche), le mécanisme de scoring (attention inter-experts vs GRU scalaire), et l'initialisation (a priori vs émergente). Les propriétés théoriques établies constituent une base formelle pour une validation empirique sur des modèles de taille intermédiaire (1B--7B paramètres), en particulier sur la question ouverte de l'estimation des experts non actifs.

\medskip
\noindent\textbf{Validation et raffinement (cette version).} La section~\ref{sec:empirique} valide empiriquement que la commutation inter-blocs est nécessaire à la composition, mais montre que des experts liés à une profondeur fixe ne généralisent pas aux compositions inédites. Le raffinement proposé --- une boucle d'experts \emph{partagés} à ré-ancrage par tokens, complétée d'un \emph{routage autonome} par démonstrations few-shot et critique de suffisance --- résout la généralisation compositionnelle zéro-shot sur le banc synthétique (précision de $1.00$ sur paires et triples inédits) et trace une voie de passage à l'échelle (c\oe ur $\to$ backbone, experts $\to$ adaptateurs, ré-ancrage $\to$ chaîne-de-pensée, critique exact $\to$ vérificateur appris, recherche $\to$ proposeur guidé). Nous séparons explicitement le \emph{mécanisme}, validé en jouet, de l'\emph{intelligence du juge} de suffisance, attendue à l'échelle LLM.

\bibliographystyle{plain}
\begin{thebibliography}{99}

\bibitem{shazeer2017outrageously}
N. Shazeer et al.
\newblock Outrageously large neural networks: The sparsely-gated mixture-of-experts layer.
\newblock \textit{ICLR}, 2017.

\bibitem{fedus2022switch}
W. Fedus, B. Zoph, N. Shazeer.
\newblock Switch transformers: Scaling to trillion parameter models with simple and efficient sparsity.
\newblock \textit{JMLR}, 2022.

\bibitem{jiang2024mixtral}
A.~Q. Jiang et al.
\newblock Mixtral of experts.
\newblock \textit{arXiv:2401.04088}, 2024.

\bibitem{qiu2025rmoe}
Z. Qiu, Z. Huang, S. Cheng, Y. Zhou, Z. Wang, I. Titov, J. Fu.
\newblock Layerwise recurrent router for mixture-of-experts.
\newblock \textit{ICLR}, 2025. arXiv:2408.06793.

\bibitem{li2022btm}
M. Li et al.
\newblock Branch-train-merge: Embarrassingly parallel training of expert language models.
\newblock \textit{arXiv:2208.03306}, 2022.

\bibitem{upit2025}
[Auteurs].
\newblock Upcycling instruction tuning from dense to mixture-of-experts.
\newblock \textit{ACL}, 2025.

\bibitem{edl2026}
[Auteurs].
\newblock Expert divergence learning for MoE-based language models.
\newblock \textit{arXiv:2603.00054}, 2026.

\bibitem{raposo2024mod}
D. Raposo et al.
\newblock Mixture-of-depths: Dynamically allocating compute in transformer-based language models.
\newblock \textit{arXiv:2404.02258}, 2024.

\bibitem{dehghani2019universal}
M. Dehghani et al.
\newblock Universal transformers.
\newblock \textit{ICLR}, 2019.

\bibitem{bengio2013estimating}
Y. Bengio, N. Léonard, A. Courville.
\newblock Estimating or propagating gradients through stochastic neurons.
\newblock \textit{arXiv:1308.3432}, 2013.

\bibitem{jang2017categorical}
E. Jang, S. Gu, B. Poole.
\newblock Categorical reparameterization with gumbel-softmax.
\newblock \textit{ICLR}, 2017.

\bibitem{dehaene2005neural}
S. Dehaene, L. Cohen.
\newblock The unique role of the visual word form area in reading.
\newblock \textit{Trends in Cognitive Sciences}, 2005.

\bibitem{wei2022cot}
J. Wei, X. Wang, D. Schuurmans, et al.
\newblock Chain-of-thought prompting elicits reasoning in large language models.
\newblock \textit{NeurIPS}, 2022.

\bibitem{banino2021ponder}
A. Banino, J. Balaguer, C. Blundell.
\newblock PonderNet: Learning to ponder.
\newblock \textit{arXiv:2107.05407}, 2021.

\bibitem{cobbe2021verifiers}
K. Cobbe, V. Kosaraju, M. Bavarian, et al.
\newblock Training verifiers to solve math word problems.
\newblock \textit{arXiv:2110.14168}, 2021.

\end{thebibliography}

\end{document}
